% !TeX program = lualatex
% ============================================================
% chapter1.tex — ch01 唯一內容源（2026-08-09 起；LaTeX 統一 U1）
% 沿革：升格自 dist/ch01/chapter1.tex（原 make_dist.py 自 fragment 轉換的成品；
%   fragment 樹已凍結，改課文一律改本檔。拍板見 ../../KICKOFF-latex-unification.md）
%   樣式源＝../../template/calcbook.sty（M-B1 拍板模板）
% 編譯（本資料夾為 CWD；aux 進 build/、成品 PDF 移入 ../../dist/ch01/）：
%   latexmk -lualatex -auxdir=../../build/aux-ch01 chapter1.tex
% ============================================================
\documentclass[a4paper,12pt,oneside]{memoir}
\makeatletter\def\input@path{{../../template/}}\makeatother
\def\cbfontsdir{../../template/fonts/inter/}
\usepackage{calcbook}
% 圖：emitter 射 `<ch>/<stem>`（見 convert.py），故 graphicspath 指向 chapters/ 上層。
% 模板內的 {../figs/} 是 chapters/<ch>/ 為 CWD 時的路徑，dist/<ch>/ 需這一行覆寫。
\graphicspath{{../../chapters/}}
\begin{document}
\chapteropener{Chapter 1}{Inverse Functions and Limits}
\begin{lead}
This chapter develops two ideas that together form the starting point of calculus. The first is the inverse of a function: when a function sends each input to a distinct output, we can ask how to recover the input from the output. This leads to one-to-one functions, inverse functions, and the inverse trigonometric functions. The second is the limit of a function. Many quantities of interest in calculus — the instantaneous rate of change of a moving object, the slope of a tangent line, the behavior of a function near a hole in its graph — depend not on the value of the function at a single point but on its behavior as the input approaches that point. We study limits first informally through tables and graphs, then through algebraic laws that let us compute them systematically, and finally through the precise \(\varepsilon\)-\(\delta\) definition that places every preceding argument on a rigorous footing.
\end{lead}
\parahead{By the end of this chapter, you will be able to:}
\begin{objectives}
  \item decide when a function is one-to-one and find its inverse when it exists;
  \item work with the inverse trigonometric functions \(\arcsin\), \(\arccos\), \(\arctan\), and their companions, together with their principal intervals;
  \item estimate the value of a limit from tables, graphs, and algebraic simplification, and handle one-sided and infinite limits;
  \item compute limits using the limit laws, factoring and rationalization, and the squeeze theorem;
  \item state and apply the precise \(\varepsilon\)-\(\delta\) definition of a limit to verify a given limit rigorously.
\end{objectives}
\sechead{1.1}{Inverse Functions}
Recall that a function \(f\) assigns a value to every input in its domain. This section studies the reverse process of recovering the input from the output. We proceed in three stages: identify which functions admit inversion (the one-to-one property), define the inverse function for those that do, and give an explicit procedure for constructing it. Before the formal definitions, two simple examples show when inversion is possible and when it is not.

\begin{warmuplist}
  \item \(f(x) = x\) on \([0, 1]\). Here \(f(0) = 0\), \(f(\tfrac{1}{2}) = \tfrac{1}{2}\), \(f(\tfrac{1}{8}) = \tfrac{1}{8}\); each input produces its own distinct output.
  \item \(f(x) = x^{2}\) on \([-1, 1]\). Here \(f(0) = 0\), \(f(\tfrac{1}{2}) = \tfrac{1}{4}\), and \(f(-\tfrac{1}{2}) = \tfrac{1}{4}\); two different inputs produce the same output.
\end{warmuplist}
The difference between these two cases is the central point. In the first, every output corresponds to exactly one input. In the second, the output \(\tfrac{1}{4}\) corresponds to both \(\tfrac{1}{2}\) and \(-\tfrac{1}{2}\), so there is no unambiguous way to go backwards from output to input.

\subsechead{One-to-one functions}
A function can be reversed only when each output comes from exactly one input. This motivates the following definition.

\begin{envdefinition}{Definition}{1.1}{}
Let \(f\) be a function with domain \(A\). We say that \(f\) is \emph{one-to-one} if \[ f(x_{1}) \neq f(x_{2}) \qquad \text{whenever } x_{1} \neq x_{2}. \] Equivalently, \(f\) is one-to-one if \[ f(x_{1}) = f(x_{2}) \implies x_{1} = x_{2}. \]

\begin{informal}
Informally, a function is one-to-one when different inputs always produce different outputs.
\end{informal}
\end{envdefinition}
\begin{workedexample}
\begin{envexample}{Example}{1.1}{}
A class of students is assigned identification numbers, arranged so that no two students share a number. Define \(f(\mathrm{student})\) to be that student’s identification number. Is \(f\) one-to-one?
\end{envexample}
\begin{envsolution}{Solution}{}{}
By construction, no two students share a number, so distinct students always produce distinct identification numbers. Therefore \(f\) is one-to-one.
\end{envsolution}
\end{workedexample}
Now we return to the two mathematical examples from the opening.

\begin{workedexample}
\begin{envexample}{Example}{1.2}{}
Determine which of the following functions are one-to-one on their given domains: \[ \begin{aligned} f(x) &= x, & x &\in [0, 1],\\ g(x) &= x^{2}, & x &\in [-1, 1]. \end{aligned} \]
\end{envexample}
\begin{envsolution}{Solution}{}{}
The function \(f(x) = x\) is one-to-one on \([0, 1]\) because each input is sent to itself, so distinct inputs always give distinct outputs.

The function \(g(x) = x^{2}\) is not one-to-one on \([-1, 1]\) because \[ \begin{aligned} g\!\left(\tfrac{1}{2}\right) &= \left(\tfrac{1}{2}\right)^{2} = \tfrac{1}{4},\\ g\!\left(-\tfrac{1}{2}\right) &= \left(-\tfrac{1}{2}\right)^{2} = \tfrac{1}{4}. \end{aligned} \] Two different inputs produce the same output, so \(g\) is not one-to-one.
\end{envsolution}
\end{workedexample}
\begin{workedexample}
\begin{envexample}{Example}{1.3}{}
Determine whether each function is one-to-one on \(\mathbb{R}\): \[ h(x) = x^{3} - 4x, \qquad k(x) = x^{3} + 4. \]
\end{envexample}
\begin{envsolution}{Solution}{}{}
For \(h\), we look for two different inputs with the same output. Since \(h(x) = x(x^{2} - 4) = x(x - 2)(x + 2)\), we have \(h(0) = h(2) = 0\) with \(0 \neq 2\), so \(h\) is \emph{not} one-to-one, even though its formula looks similar to \(x^{3}\).

For \(k\), suppose \(k(x_{1}) = k(x_{2})\). Then \(x_{1}^{3} + 4 = x_{2}^{3} + 4\), so \(x_{1}^{3} = x_{2}^{3}\), and taking cube roots gives \(x_{1} = x_{2}\). Therefore \(k\) \emph{is} one-to-one. The lesson: two formulas of the same general shape can behave differently; when in doubt, test the definition directly.
\end{envsolution}
\end{workedexample}
Looking at the graph of a function gives a quick way to decide whether it is one-to-one.

\begin{envremark}{Remark}{1.1}{Horizontal line test}
A horizontal line \(y = c\) meets the graph of \(f\) at a point \((x, c)\) exactly when \(f(x) = c\). Therefore, \emph{a function is one-to-one if and only if no horizontal line intersects its graph more than once}. If the line \(y = c\) meets the graph at two distinct points \((x_{1}, c)\) and \((x_{2}, c)\), then \(f(x_{1}) = c = f(x_{2})\) with \(x_{1} \neq x_{2}\), so \(f\) fails to be one-to-one. Conversely, if every horizontal line meets the graph at most once, each output corresponds to at most one input. Figure 1.1 shows the two situations.
\end{envremark}
\begin{figureblock}
  \includegraphics[width=72.35mm]{ch01/figs/hlt-1}\hspace{4.3mm}\includegraphics[width=72.35mm]{ch01/figs/hlt-2}
\figcaption{Figure 1.1}{Geometric interpretation of the horizontal line test.}
\end{figureblock}
\pagebreakbefore
\subsechead{Inverse functions}
If \(f\) is one-to-one, then we can define a second function that undoes \(f\): given an output of \(f\), it returns the unique input that produced it.

The idea is simply that the inverse trades the roles of input and output. For \(f\), the question is “given the input \(x\), what is the output?” The inverse reads the same correspondence in the other direction: “given that output, which input produced it?” This is why the domain and range change places: the outputs of \(f\) are exactly the inputs of \(f^{-1}\), and the inputs of \(f\) are exactly the outputs of \(f^{-1}\).

\begin{envdefinition}{Definition}{1.2}{}
Let \(f\) be a one-to-one function with domain \(A\) and range \(B\). The \emph{inverse function} of \(f\), denoted by \(f^{-1}\), has domain \(B\) and range \(A\), and is defined, for every \(y \in B\), by \[ f^{-1}(y) = x \quad \Longleftrightarrow \quad f(x) = y. \]
\end{envdefinition}
\begin{envcaution}{Caution}{1.1}{}
The superscript \(-1\) in \(f^{-1}\) denotes the \emph{inverse function}, not a reciprocal. That is, \(f^{-1}(x)\) is the inverse evaluated at \(x\), whereas \(\bigl(f(x)\bigr)^{-1} = 1 / f(x)\) is the reciprocal of the value \(f(x)\); these are generally different. The same warning applies to the inverse trigonometric notation \(\sin^{-1} x\) in the next section.
\end{envcaution}
\begin{envremark}{Remark}{1.2}{}
When we focus on the inverse function itself, it is often convenient to use \(x\) as its independent variable. In that case, we write \[ f^{-1}(x) = y \quad \Longleftrightarrow \quad f(y) = x. \]
\end{envremark}
\pagebreakbefore
\begin{envtheorem}{Theorem}{1.1}{}
A function has an inverse if and only if it is one-to-one.
\end{envtheorem}
\begin{envproof}{Proof}{}{}
Suppose first that \(f\) has an inverse \(f^{-1}\). If \(f(x_{1}) = f(x_{2})\), applying \(f^{-1}\) to both sides gives \[ f^{-1}\!\bigl(f(x_{1})\bigr) = f^{-1}\!\bigl(f(x_{2})\bigr), \] so \(x_{1} = x_{2}\). Therefore \(f\) is one-to-one. Conversely, suppose \(f\) is one-to-one with domain \(A\) and range \(B\). For each \(y \in B\) there exists some \(x \in A\) with \(f(x) = y\), and the one-to-one property forces this \(x\) to be unique. The correspondence \(y \mapsto x\) therefore defines an inverse function \(f^{-1} : B \to A\).
\end{envproof}
\begin{workedexample}
\begin{envexample}{Example}{1.4}{}
Find the inverse of each function.

\begin{enumerate}
  \item \(f(x) = x\).
  \item \(g(x) = x^{3}\).
\end{enumerate}
\end{envexample}
\begin{envsolution}{Solution}{}{}
\begin{enumerate}
  \item Since \(f(x) = x\) leaves every input unchanged, \(f\) is its own inverse, so \(f^{-1}(x) = x\).
  \item To find the inverse of \(g\), let \(y = x^{3}\). Solving for \(x\) gives \(x = \sqrt[3]{y}\). Interchanging the names of the variables, \(g^{-1}(x) = \sqrt[3]{x}\). As a check, \(g\!\bigl(g^{-1}(x)\bigr) = \bigl(\sqrt[3]{x}\bigr)^{3} = x\).
\end{enumerate}
\end{envsolution}
\end{workedexample}
The defining relation \[ f^{-1}(y) = x \Longleftrightarrow f(x) = y \] yields two composition identities.

\begin{envproposition}{Proposition}{1.1}{}
Let \(f\) be a one-to-one function with domain \(A\) and range \(B\). Then \[ \begin{aligned} f^{-1}\!\bigl(f(x)\bigr) &= x \quad \text{for every } x \in A, \\ f\!\bigl(f^{-1}(y)\bigr) &= y \quad \text{for every } y \in B. \end{aligned} \]
\end{envproposition}
Figure 1.2 illustrates both identities: \(f\) and \(f^{-1}\) travel between \(A\) and \(B\) in opposite directions, and composing them recovers the original input.

\begin{figureblock}
  \includegraphics[width=90.49mm]{ch01/figs/fig-map}
\figcaption{Figure 1.2}{A one-to-one function and its inverse map between \(A\) and \(B\).}
\end{figureblock}
Graphically, if the graph of \(f^{-1}\) is drawn in the same coordinate plane as the graph of \(f\), the two graphs are reflections of one another across the line \(y = x\). The point \((a, b)\) lies on the graph of \(f\) exactly when \((b, a)\) lies on the graph of \(f^{-1}\). You will use this machinery again in Chapter 4, where the natural logarithm is defined as the inverse of the exponential function. Its graph is exactly this kind of reflection.

\begin{workedexample}
\begin{envexample}{Example}{1.5}{}
The function \(f(x) = x^{2}\) is not one-to-one on \(\mathbb{R}\). Restrict its domain so that the restricted function is one-to-one, and find the inverse of the restricted function.
\end{envexample}
\begin{envsolution}{Solution}{}{}
This example illustrates why not every function can be reversed, and how restricting the domain repairs it. On the interval \([0, \infty)\) the function \(f(x) = x^{2}\) is one-to-one: if \(x_{1}^{2} = x_{2}^{2}\) with \(x_{1}, x_{2} \ge 0\), then \(x_{1} = x_{2}\). To find the inverse, apply the same procedure: from \(y = x^{2}\) with \(x \ge 0\), solving for \(x\) gives the positive root \(x = \sqrt{y}\), since \(x \ge 0\). Interchanging the names of the variables gives \[ f^{-1}(x) = \sqrt{x}, \qquad x \ge 0. \] This is exactly why \(\sqrt{x}\) means the \emph{positive} square root: the convention makes the square root the inverse of the restricted squaring function. Figure 1.3 shows the restricted parabola, its inverse, and the mirror line \(y = x\). Restricting first and then inverting is the key to the inverse trigonometric functions of the next section.

\begin{figureblock}
  \includegraphics[width=87.38mm]{ch01/figs/restrict-x2}
\figcaption{Figure 1.3}{Restricted squaring, its inverse \(\sqrt{x}\), and the mirror line \(y = x\).}
\end{figureblock}
\end{envsolution}
\end{workedexample}
\subsechead{Finding the inverse of a function}
When \(f\) is one-to-one, we can usually construct its inverse by a short procedure.

\begin{envstrategy}{Strategy}{1.1}{Finding the inverse of a one-to-one function}
\begin{steps}
  \item Write \(y = f(x)\).
  \item Solve this equation for \(x\) in terms of \(y\).
  \item Interchange the names \(x\) and \(y\). The resulting equation gives \(y = f^{-1}(x)\).
\end{steps}
\end{envstrategy}
Note that this procedure assumes \(f\) is already one-to-one: applying it blindly to \(y = x^{2}\), for instance, yields \(x = \pm\sqrt{y}\): two values, not a function. Always verify one-to-one first, or restrict the domain (as in Example 1.5).

\begin{workedexample}
\begin{envexample}{Example}{1.6}{}
Find the inverse of \(f(x) = x^{3} + 2\).
\end{envexample}
\begin{envsolution}{Solution}{}{}
Let \(y = x^{3} + 2\). Solving for \(x\) gives \[ \begin{aligned} x^{3} &= y - 2,\\ x &= \sqrt[3]{y - 2}. \end{aligned} \] Interchanging the variables, \(y = \sqrt[3]{x - 2}\), so \(f^{-1}(x) = \sqrt[3]{x - 2}\). As a check, \[ f\!\bigl(f^{-1}(x)\bigr) = \left(\sqrt[3]{x - 2}\right)^{3} + 2 = (x - 2) + 2 = x. \]
\end{envsolution}
\end{workedexample}
\begin{workedexample}
\begin{envexample}{Example}{1.7}{}
Find the inverse of \(f(x) = \frac{3}{x - 5}\) (for \(x \neq 5\)), and verify your answer by computing both \(f\!\bigl(f^{-1}(x)\bigr)\) and \(f^{-1}\!\bigl(f(x)\bigr)\).
\end{envexample}
\begin{envsolution}{Solution}{}{}
Let \(y = \frac{3}{x - 5}\). Solving for \(x\): \[ x - 5 = \frac{3}{y}, \quad\text{so}\quad x = 5 + \frac{3}{y}. \] Interchanging the variables, \[ f^{-1}(x) = 5 + \frac{3}{x} = \frac{5x + 3}{x}, \qquad x \neq 0. \] Verification, in both directions: \[ \begin{aligned} f\!\bigl(f^{-1}(x)\bigr) &= \frac{3}{\bigl(5 + \tfrac{3}{x}\bigr) - 5} = \frac{3}{3/x} = x, \\ f^{-1}\!\bigl(f(x)\bigr) &= 5 + \frac{3}{3/(x - 5)} = 5 + (x - 5) = x. \end{aligned} \] Both compositions return the input, confirming Proposition 1.1. Note how the domains pair up: \(f\) excludes \(x = 5\) and \(f^{-1}\) excludes \(x = 0\), matching the fact that \(f\) never takes the value \(0\) and \(f^{-1}\) never takes the value \(5\).
\end{envsolution}
\end{workedexample}
\begin{workedexample}
\begin{envexample}{Example}{1.8}{}
The function \(f(t) = \tfrac{9}{5} t + 32\) converts a temperature of \(t\) degrees Celsius into degrees Fahrenheit. What practical question does \(f^{-1}\) answer? Find a formula for \(f^{-1}(t)\), and verify your formula by computing \(f^{-1}\!\bigl(f(t)\bigr)\).
\end{envexample}
\begin{envsolution}{Solution}{}{}
Since \(f\) converts Celsius to Fahrenheit, its inverse \(f^{-1}\) converts a Fahrenheit reading back into Celsius. So \(f^{-1}(t)\) is the Celsius value whose Fahrenheit reading is \(t\); applying \(f\) to it returns \(t\), giving \(t = \tfrac{9}{5} f^{-1}(t) + 32\). Solving, \[ f^{-1}(t) = \tfrac{5}{9}\,(t - 32). \] As a check, \[ f^{-1}\!\bigl(f(t)\bigr) = \tfrac{5}{9}\!\left(\tfrac{9}{5} t + 32 - 32\right) = \tfrac{5}{9}\cdot\tfrac{9}{5}\,t = t. \]
\end{envsolution}
\end{workedexample}

\sechead{1.2}{Inverse Trigonometric Functions}
The trigonometric functions are not one-to-one on all of \(\mathbb{R}\), so they do not have inverse functions on their full domains. Inverse trigonometric functions arise whenever we know a ratio and want to recover an angle: from the slope of a ramp, we get an angle of inclination via \(\arctan\); from a ratio of sides in a right triangle, we get an angle via \(\arcsin\) or \(\arccos\). To build these inverses, we must first restrict the domain of each trigonometric function to an interval on which it is one-to-one. You will use these functions again in Chapter 3, where the derivatives of \(\arcsin\), \(\arccos\), and \(\arctan\) are computed. The principal ranges fixed in this section are exactly what determine those formulas.

Consider the sine function on \(\mathbb{R}\). As Figure 1.4 shows, a horizontal line can intersect the graph of \(y = \sin x\) more than once, so \(\sin\) is not one-to-one on \(\mathbb{R}\).

\begin{figureblock}
  \includegraphics[width=132.82mm]{ch01/figs/sine-not-1to1}
\figcaption{Figure 1.4}{The sine function on \(\mathbb{R}\) is not one-to-one.}
\end{figureblock}
\subsechead{The inverse sine function}
To define the inverse sine, we restrict \(\sin x\) to the interval \[ -\tfrac{\pi}{2} \le x \le \tfrac{\pi}{2}. \] On this interval, \(\sin x\) is one-to-one and its values cover the range \([-1, 1]\). The restricted branch is shown in Figure 1.5.

\begin{figureblock}
  \includegraphics[width=90.49mm]{ch01/figs/restricted-sine}
\figcaption{Figure 1.5}{Restricted branch of the sine function used to define \(\arcsin x\).}
\end{figureblock}
\begin{envdefinition}{Definition}{1.3}{}
The \emph{inverse sine function}, denoted \(\arcsin x\) or \(\sin^{-1} x\), is defined by \[ \arcsin x = y \quad \Longleftrightarrow \quad \sin y = x \ \text{ and } \ -\tfrac{\pi}{2} \le y \le \tfrac{\pi}{2}. \] Its domain is \([-1, 1]\) and its range is \(\left[-\tfrac{\pi}{2}, \tfrac{\pi}{2}\right]\).
\end{envdefinition}
\begin{envcaution}{Caution}{1.2}{}
The notation \(\sin^{-1} x\) is also used for \(\arcsin x\). It does \emph{not} mean the reciprocal \(\bigl(\sin x\bigr)^{-1} = 1 / \sin x\). When reading a source that uses \(\sin^{-1}\), check the context to decide which is meant.
\end{envcaution}
The composition identities for \(\arcsin\) take the following form.

\begin{envproposition}{Proposition}{1.2}{Inverse sine identities}
\[ \begin{aligned} \sin(\arcsin x) = x, &\qquad -1 \le x \le 1,\\ \arcsin(\sin x) = x, &\qquad -\tfrac{\pi}{2} \le x \le \tfrac{\pi}{2}. \end{aligned} \]
\end{envproposition}
\begin{envcaution}{Caution}{1.3}{}
The identity \(\arcsin(\sin x) = x\) holds only when \(x \in \left[-\tfrac{\pi}{2}, \tfrac{\pi}{2}\right]\). Outside this interval, the right-hand side is replaced by the unique value in \(\left[-\tfrac{\pi}{2}, \tfrac{\pi}{2}\right]\) with the same sine; for instance, \[ \arcsin(\sin \pi) = \arcsin(0) = 0 \neq \pi. \]
\end{envcaution}
\begin{envstrategy}{Strategy}{1.2}{The reference-triangle method}
To evaluate a trigonometric function of an inverse trigonometric expression, such as \(\tan(\arcsin x)\) or \(\cos(\arctan x)\):

\begin{steps}
  \item Set \(\theta\) equal to the inverse trigonometric expression, so that the known ratio is read directly from the definition (e.g. \(\theta = \arcsin x\) gives \(\sin\theta = x\)).
  \item Draw a right triangle whose side lengths are the magnitudes (absolute values) of the known ratio. Keep the sign separate, and use the principal-value interval to settle it. That interval is the range fixed in the inverse function’s definition.
  \item Find the missing side by the Pythagorean theorem.
  \item Read the desired value from the completed triangle. Determine the sign from the principal-value interval, not from the triangle alone.
\end{steps}
\end{envstrategy}
\begin{workedexample}
\begin{envexample}{Example}{1.9}{}
Evaluate the following expressions.

\begin{enumerate}
  \item \(\arcsin\!\left(\tfrac{1}{2}\right)\).
  \item \(\tan\bigl(\arcsin\!\left(\tfrac{1}{3}\right)\bigr)\).
\end{enumerate}
\end{envexample}
\begin{envsolution}{Solution}{}{}
\begin{enumerate}
  \item We need the angle \(y \in \left[-\tfrac{\pi}{2}, \tfrac{\pi}{2}\right]\) with \(\sin y = \tfrac{1}{2}\). Since \(\sin(\pi/6) = \tfrac{1}{2}\), we have \(\arcsin\!\left(\tfrac{1}{2}\right) = \tfrac{\pi}{6}\).
  \item Let \(\theta = \arcsin\!\left(\tfrac{1}{3}\right)\), so \(\sin\theta = \tfrac{1}{3}\) with \(\theta \in \left[-\tfrac{\pi}{2}, \tfrac{\pi}{2}\right]\). Because \(\sin\theta > 0\), the angle \(\theta\) lies in the first quadrant, and in a right triangle with hypotenuse \(3\) and opposite side \(1\), the adjacent side is \[ a = \sqrt{3^{2} - 1^{2}} = \sqrt{8} = 2\sqrt{2}. \]
\begin{figureblock}
  \includegraphics[width=53.45mm]{ch01/figs/arcsin-triangle}
\figcaption{Figure 1.6}{Right triangle with \(\sin\theta = \tfrac{1}{3}\).}
\end{figureblock}
Therefore \(\tan\bigl(\arcsin\!\left(\tfrac{1}{3}\right)\bigr) = \tan\theta = \tfrac{1}{2\sqrt{2}}\).
\end{enumerate}
\end{envsolution}
\end{workedexample}
\subsechead{The inverse cosine function}
As with sine, we first choose an interval on which the function is one-to-one and still runs through all of its values. Cosine decreases steadily from \(1\) to \(-1\) as the angle goes from \(0\) to \(\pi\), so the interval \([0, \pi]\) sweeps the entire range \([-1, 1]\) exactly once, the natural analogue of the choice \(\left[-\tfrac{\pi}{2}, \tfrac{\pi}{2}\right]\) made for sine.

The cosine function is one-to-one on the interval \(0 \le x \le \pi\), shown in Figure 1.7.

\begin{figureblock}
  \includegraphics[width=90.49mm]{ch01/figs/restricted-cosine}
\figcaption{Figure 1.7}{Restricted branch of the cosine function used to define \(\arccos x\).}
\end{figureblock}
\begin{envdefinition}{Definition}{1.4}{}
The \emph{inverse cosine function}, denoted \(\arccos x\) or \(\cos^{-1} x\), is defined by \[ \arccos x = y \quad \Longleftrightarrow \quad \cos y = x \ \text{ and } \ 0 \le y \le \pi. \] Its domain is \([-1, 1]\) and its range is \([0, \pi]\).
\end{envdefinition}
\begin{envproposition}{Proposition}{1.3}{Inverse cosine identities}
\[ \begin{aligned} \cos(\arccos x) = x, &\qquad -1 \le x \le 1,\\ \arccos(\cos x) = x, &\qquad 0 \le x \le \pi. \end{aligned} \]
\end{envproposition}
\begin{workedexample}
\begin{envexample}{Example}{1.10}{}
Find the domain of each function.

\begin{enumerate}
  \item \(f(x) = \arcsin(\cos x)\).
  \item \(h(x) = \sin(\arccos x)\).
\end{enumerate}
\end{envexample}
\begin{envsolution}{Solution}{}{}
\begin{enumerate}
  \item Whatever \(x\) is, \(\cos x\) lies in \([-1, 1]\), and \([-1, 1]\) is exactly the domain of \(\arcsin\). So every real \(x\) is allowed: the domain of \(f\) is \(\mathbb{R}\).
  \item The inner function \(\arccos x\) requires \(x \in [-1, 1]\); its output is an angle, and \(\sin\) accepts any angle. So the domain of \(h\) is \([-1, 1]\).
\end{enumerate}
To find the domain of a composition, first make sure the input is allowed by the \emph{inner} function. Then make sure the output of the inner function is allowed by the \emph{outer} function.
\end{envsolution}
\end{workedexample}
\begin{workedexample}
\begin{envexample}{Example}{1.11}{}
A particle starts moving at time \(t = 10\), and it bobs up and down so that its height at time \(t \ge 10\) is \(\cos t\). True or false: the particle has height \(1\) at time \(t = \arccos(1)\).
\end{envexample}
\begin{envsolution}{Solution}{}{}
False. The equation \(\cos t = 1\) holds for infinitely many \(t\), but \(\arccos\) returns only the principal value \(t = 0\). The particle is not even moving at \(t = 0\), since its motion starts at \(t = 10\). The particle does reach height \(1\) at the times \(t = 2\pi n\) for integer \(n\); because the motion exists only for \(t \ge 10\), this requires \(2\pi n \ge 10\), i.e. \(n \ge 2\), but \(\arccos(1) = 0\) is not one of them.
\end{envsolution}
\end{workedexample}
\subsechead{The inverse tangent function}
On the open interval \(\left(-\tfrac{\pi}{2}, \tfrac{\pi}{2}\right)\), the tangent function climbs continuously from \(-\infty\) to \(+\infty\), taking each real number as a value exactly once. This is why the inverse tangent, unlike \(\arcsin\) and \(\arccos\), accepts \emph{every} real number as input: its domain is all of \(\mathbb{R}\) (Figure 1.9).

The tangent function is one-to-one on the interval \(-\tfrac{\pi}{2} < x < \tfrac{\pi}{2}\), shown in Figure 1.8.

\begin{figureblock}
  \includegraphics[width=79.9mm]{ch01/figs/restricted-tangent}
\figcaption{Figure 1.8}{Restricted branch of the tangent function used to define \(\arctan x\).}
\end{figureblock}
\begin{envdefinition}{Definition}{1.5}{}
The \emph{inverse tangent function}, denoted \(\arctan x\) or \(\tan^{-1} x\), is defined by \[ \arctan x = y \quad \Longleftrightarrow \quad \tan y = x \ \text{ and } \ -\tfrac{\pi}{2} < y < \tfrac{\pi}{2}. \] Its domain is \(\mathbb{R}\) and its range is \(\left(-\tfrac{\pi}{2}, \tfrac{\pi}{2}\right)\).
\end{envdefinition}
Figure 1.9 shows this reflection. In the reflected graph, the vertical asymptotes of the restricted tangent branch become horizontal asymptotes of the inverse tangent graph.

\begin{figureblock}
  \includegraphics[width=79.9mm]{ch01/figs/arctan-reflection}
\figcaption{Figure 1.9}{The restricted tangent branch and its inverse, reflected across \(y = x\).}
\end{figureblock}
\begin{envproposition}{Proposition}{1.4}{Inverse tangent identities}
\[ \begin{aligned} \tan(\arctan x) = x, &\qquad x \in \mathbb{R},\\ \arctan(\tan x) = x, &\qquad -\tfrac{\pi}{2} < x < \tfrac{\pi}{2}. \end{aligned} \]
\end{envproposition}
\begin{envcaution}{Caution}{1.4}{}
As with the sine case, the “cancellation” identities hold only on each principal range. That is, \(\arccos(\cos x) = x\) only when \(x \in [0, \pi]\), and \(\arctan(\tan x) = x\) only when \(x \in \left(-\tfrac{\pi}{2}, \tfrac{\pi}{2}\right)\). Outside these intervals the left-hand side returns the equivalent angle that \emph{does} lie in the principal range; for instance, \[ \begin{aligned} \arccos(\cos 2\pi) &= \arccos(1) = 0 \neq 2\pi, \\ \arctan(\tan \pi) &= \arctan(0) = 0 \neq \pi. \end{aligned} \] By contrast, the reverse direction is always safe on the natural domain: \(\cos(\arccos y) = y\) for every \(y \in [-1, 1]\), and \(\tan(\arctan y) = y\) for every real \(y\).
\end{envcaution}
The next two examples involve the \emph{reciprocal} trigonometric functions secant, cosecant, and cotangent. If your precalculus course made little use of them, their definitions are all you need: for any angle \(\theta\) at which the denominator is nonzero, \[ \sec\theta = \frac{1}{\cos\theta}, \qquad \csc\theta = \frac{1}{\sin\theta}, \qquad \cot\theta = \frac{\cos\theta}{\sin\theta}. \] Since \(\cos\theta\) is the ratio adjacent over hypotenuse in a right triangle, \(\sec\theta\) is the ratio hypotenuse over adjacent, and similarly for the other two.

Dividing the Pythagorean identity \(\sin^{2}\theta + \cos^{2}\theta = 1\) through by \(\cos^{2}\theta\) gives the companion identity \(1 + \tan^{2}\theta = \sec^{2}\theta\), which Example 1.13 uses. Note the contrast with Caution 1.2: the superscript in \(\sin^{-1} x\) does \emph{not} denote a reciprocal, but \(\sec\) and \(\csc\) really are reciprocals (and \(\cot\) is the inverted ratio \(\cos/\sin\)). Appendix A.1 develops these three functions more fully.

\begin{workedexample}
\begin{envexample}{Example}{1.12}{}
Find the exact value of each expression.

\begin{enumerate}
  \item \(\sin\bigl(\arccos(-\tfrac{1}{2})\bigr)\).
  \item \(\sec(\arctan 10)\).
  \item \(\tan\bigl(\arcsin(-\tfrac{2\sqrt{5}}{5})\bigr)\).
\end{enumerate}
\end{envexample}
\begin{envsolution}{Solution}{}{}
\begin{enumerate}
  \item \(\arccos(-\tfrac{1}{2})\) is the angle in \([0, \pi]\) whose cosine is \(-\tfrac{1}{2}\): that is \(\tfrac{2\pi}{3}\), and \(\sin\tfrac{2\pi}{3} = \tfrac{\sqrt{3}}{2}\). Note the principal value lands in the \emph{second} quadrant. Negative inputs to \(\arccos\) always give a principal value greater than \(\tfrac{\pi}{2}\).
  \item Let \(\theta = \arctan 10\), so \(\tan\theta = 10\) with \(\theta \in \left(0, \tfrac{\pi}{2}\right)\). In a right triangle with opposite side \(10\) and adjacent side \(1\), the hypotenuse is \(\sqrt{101}\), so \(\sec(\arctan 10) = \sqrt{101}\).
\begin{figureblock}
  \includegraphics[width=53.45mm]{ch01/figs/arctan10-triangle}
\figcaption{Figure 1.10}{Right triangle with \(\tan\theta = 10\).}
\end{figureblock}

  \item Let \(\theta = \arcsin\bigl(-\tfrac{2\sqrt{5}}{5}\bigr)\), so \(\sin\theta = -\tfrac{2}{\sqrt{5}}\) with \(\theta \in \left[-\tfrac{\pi}{2}, \tfrac{\pi}{2}\right]\). Because \(\sin\theta < 0\), the angle lies in the fourth quadrant, where cosine is positive: \(\cos\theta = \sqrt{1 - \tfrac{4}{5}} = \tfrac{1}{\sqrt{5}}\). Therefore \(\tan\theta = \frac{-2/\sqrt{5}}{1/\sqrt{5}} = -2\). The sign came from the \emph{range} of \(\arcsin\), not from guessing.
\end{enumerate}
\end{envsolution}
\end{workedexample}
\begin{workedexample}
\begin{envexample}{Example}{1.13}{}
Simplify the expression \(\cos(\arctan x)\).
\end{envexample}
\begin{envsolution}{Solution}{}{}
Let \(y = \arctan x\), so \(\tan y = x\) with \(-\tfrac{\pi}{2} < y < \tfrac{\pi}{2}\). We want \(\cos y\). Using the identity \(1 + \tan^{2} y = \sec^{2} y\), \[ \sec^{2} y = 1 + x^{2}. \] Since \(y \in \left(-\tfrac{\pi}{2}, \tfrac{\pi}{2}\right)\) we have \(\cos y > 0\) and hence \(\sec y > 0\), so \(\sec y = \sqrt{1 + x^{2}}\). Therefore \[ \cos(\arctan x) = \cos y = \frac{1}{\sec y} = \frac{1}{\sqrt{1 + x^{2}}}. \]

\begin{figureblock}
  \includegraphics[width=64.03mm]{ch01/figs/arctan-general-triangle}
\figcaption{Figure 1.11}{Right triangle with \(\tan y = x\).}
\end{figureblock}
\end{envsolution}
\end{workedexample}
\subsechead{The remaining inverse trigonometric functions}
The remaining inverse trigonometric functions are used less frequently, but they follow the same pattern: restrict each to a one-to-one branch, then invert. Each definition below specifies a \emph{principal range}: the interval its output angle is restricted to. The corresponding input domains are collected together after all three.

\begin{envdefinition}{Definition}{1.6}{}
The \emph{inverse cosecant function}, denoted \(\arccsc x\) or \(\csc^{-1} x\), is defined by \[ \arccsc x = y \quad \Longleftrightarrow \quad \csc y = x \ \text{ and } \ y \in \left(0, \tfrac{\pi}{2}\right] \cup \left(\pi, \tfrac{3\pi}{2}\right]. \]
\end{envdefinition}
\begin{envdefinition}{Definition}{1.7}{}
The \emph{inverse secant function}, denoted \(\arcsec x\) or \(\sec^{-1} x\), is defined by \[ \arcsec x = y \quad \Longleftrightarrow \quad \sec y = x \ \text{ and } \ y \in \left[0, \tfrac{\pi}{2}\right) \cup \left[\pi, \tfrac{3\pi}{2}\right). \]
\end{envdefinition}
\begin{envdefinition}{Definition}{1.8}{}
The \emph{inverse cotangent function}, denoted \(\arccot x\) or \(\cot^{-1} x\), is defined by \[ \arccot x = y \quad \Longleftrightarrow \quad \cot y = x \ \text{ and } \ y \in (0, \pi). \]
\end{envdefinition}
The domains of \(\arccsc\) and \(\arcsec\) are \((-\infty, -1] \cup [1, \infty)\), while the domain of \(\arccot\) is \(\mathbb{R}\). The principal ranges are the intervals specified in the definitions above.

\begin{envremark}{Remark}{1.3}{}
For \(\arcsec x\), another common choice of principal range is \(y \in \left[0, \tfrac{\pi}{2}\right) \cup \left(\tfrac{\pi}{2}, \pi\right]\). In this text we use \(y \in \left[0, \tfrac{\pi}{2}\right) \cup \left[\pi, \tfrac{3\pi}{2}\right)\) because it makes the later derivative formula for the inverse secant simpler. When consulting other sources or software, always check which convention is in use, since different branch choices can lead to different values.
\end{envremark}
\begin{envremark}{Remark}{1.4}{}
For \(\arccsc x\), another common choice of principal range is \(y \in \left[-\tfrac{\pi}{2}, 0\right) \cup \left(0, \tfrac{\pi}{2}\right]\). In this text we use \(y \in \left(0, \tfrac{\pi}{2}\right] \cup \left(\pi, \tfrac{3\pi}{2}\right]\) so that the chosen branch parallels the one selected for \(\arcsec x\). Once the principal range is fixed consistently, either convention gives a valid inverse function.
\end{envremark}
\begin{workedexample}
\begin{envexample}{Example}{1.14}{}
Let \(f(x) = \arcsin x + \arccsc x\). Find the domain of \(f\), and evaluate \(f\) at every point of its domain.
\end{envexample}
\begin{envsolution}{Solution}{}{}
\(\arcsin\) requires \(|x| \le 1\) while \(\arccsc\) requires \(|x| \ge 1\), so the domain is just the two points \(x = \pm 1\). At \(x = 1\): \(\arcsin 1 = \tfrac{\pi}{2}\) and \(\arccsc 1 = \tfrac{\pi}{2}\) (the angle in \(\left(0, \tfrac{\pi}{2}\right]\) with cosecant \(1\)), so \(f(1) = \pi\). At \(x = -1\): \(\arcsin(-1) = -\tfrac{\pi}{2}\), and under this text’s principal range for \(\arccsc\) the angle with cosecant \(-1\) is \(\arccsc(-1) = \tfrac{3\pi}{2}\), so \(f(-1) = -\tfrac{\pi}{2} + \tfrac{3\pi}{2} = \pi\). So \(f\) is the constant \(\pi\) on its two-point domain.

Under the alternative convention of Remark 1.4 (\(\arccsc(-1) = -\tfrac{\pi}{2}\)) one gets \(f(-1) = -\pi\) instead. This is a concrete instance of branch conventions changing answers, and it is why this text fixes its conventions explicitly.
\end{envsolution}
\end{workedexample}

\sechead{1.3}{The Limit of a Function}
A frequent setting in calculus is the need to study a ratio of changes. If \(s(t)\) is the position of a moving object at time \(t\), the ratio \[ \frac{s(t + h) - s(t)}{h} \] is the \emph{average velocity} of the object from time \(t\) to time \(t + h\). Very often we want the \emph{instantaneous velocity} — how fast the object is moving at a single moment — rather than the average. To describe this, we let \(h\) shrink toward zero and ask what value the ratio approaches. More generally, for any function \(f\), the ratio \[ \frac{f(y) - f(x)}{y - x} \] measures how \(f\) changes between two inputs. We may want to know what this ratio does as \(y\) gets closer and closer to \(x\), without ever equalling \(x\). This is the motivation for the concept of a limit.

\begin{envdefinition}{Definition}{1.9}{}
Suppose that \(f(x)\) is defined for all \(x\) near a number \(a\), except possibly at \(a\) itself. We write \[ \lim_{x \to a} f(x) = L \] and say that the \emph{limit} of \(f(x)\) as \(x\) approaches \(a\) is \(L\) if \(f(x)\) can be made arbitrarily close to \(L\) whenever \(x\) is sufficiently close to \(a\) but not equal to \(a\).
\end{envdefinition}
An alternative notation for the same idea is \[ f(x) \to L \qquad \text{as } x \to a, \] read “\(f(x)\) tends to \(L\) as \(x\) tends to \(a\).”

Notice that the definition deliberately says nothing about \(f(a)\): when we consider \(\lim_{x \to a} f(x)\), we are concerned with the values of \(f(x)\) for \(x\) near \(a\), not with the value of \(f(a)\). In particular, \(f(a)\) need not be defined at all. Figure 1.12 shows three functions whose values at \(a\) are different (and one is undefined), yet which all share the same limit as \(x \to a\).

\begin{figureblock}
  \includegraphics[width=42.86mm]{ch01/figs/limit-same-near-a-1}\hspace{6mm}\includegraphics[width=42.86mm]{ch01/figs/limit-same-near-a-2}\hspace{6mm}\includegraphics[width=42.86mm]{ch01/figs/limit-same-near-a-3}
\figcaption{Figure 1.12}{Three functions with the same limit as \(x \to a\).}
\end{figureblock}
\begin{workedexample}
\begin{envexample}{Example}{1.15}{}
The graph of a function \(f\) is shown in Figure 1.13. Evaluate the following, or state that the limit does not exist:

\begin{enumerate}
  \item \(\lim_{x \to -2} f(x)\).
  \item \(\lim_{x \to 0} f(x)\).
  \item \(\lim_{x \to 2} f(x)\).
\end{enumerate}
\begin{figureblock}
  \includegraphics[width=101.07mm]{ch01/figs/read-limit-graph}
\figcaption{Figure 1.13}{A function whose limit at \(x = 2\) differs from its value there.}
\end{figureblock}
\end{envexample}
\begin{envsolution}{Solution}{}{}
\begin{enumerate}
  \item As \(x\) approaches \(-2\) from either side, the curve approaches height \(1\), so \(\lim_{x \to -2} f(x) = 1\).
  \item As \(x\) approaches \(0\), the curve approaches height \(0\), so \(\lim_{x \to 0} f(x) = 0\).
  \item As \(x\) approaches \(2\), the curve approaches height \(2\). The open dot at \((2, 2)\) and the solid dot at \((2, -2)\) tell us \(f(2) = -2\), but the limit ignores the value \emph{at} the point: \(\lim_{x \to 2} f(x) = 2\), even though \(f(2) = -2\).
\end{enumerate}
\end{envsolution}
\end{workedexample}
\begin{workedexample}
\begin{envexample}{Example}{1.16}{}
Sketch the graph of a function \(f\) with \(\lim_{x \to 3} f(x) = 10\) and \(f(3) = 0\).
\end{envexample}
\begin{envsolution}{Solution}{}{}
Many answers are possible. Take any curve that approaches height \(10\) as \(x \to 3\), for instance a straight line through \((3, 10)\). Then move the single point at \(x = 3\) down to height \(0\), marking an open dot at \((3, 10)\) and a solid dot at \((3, 0)\). Near \(x = 3\), but not at it, the values approach \(10\), so \(\lim_{x \to 3} f(x) = 10\); the relocated point makes \(f(3) = 0\). The limit and the value at the point are independent, so nothing prevents this mismatch.
\end{envsolution}
\end{workedexample}
\begin{workedexample}
\begin{envexample}{Example}{1.17}{}
Guess the value of \[ \lim_{x \to 1} \frac{x - 1}{x^{2} - 1}. \]
\end{envexample}
\begin{envsolution}{Solution}{}{}
The expression is not defined at \(x = 1\), but we can examine its values for \(x\) near \(1\):

\begin{datatable}{rcccccc}
  \(x\) & 0.9 & 0.99 & 0.999 & 1.001 & 1.01 & 1.1 \\
  \midrule
  \(\frac{x - 1}{x^{2} - 1}\) & 0.5263 & 0.5025 & 0.5003 & 0.4998 & 0.4975 & 0.4762 \\
\end{datatable}
The values are getting close to \(\tfrac{1}{2}\) as \(x\) approaches \(1\), so we guess \[ \lim_{x \to 1} \frac{x - 1}{x^{2} - 1} = \tfrac{1}{2}. \]
\end{envsolution}
\end{workedexample}
\begin{workedexample}
\begin{envexample}{Example}{1.18}{}
Guess the value of \[ \lim_{t \to 0} \frac{\sqrt{t^{2} + 9} - 3}{t^{2}}. \]
\end{envexample}
\begin{envsolution}{Solution}{}{}
We evaluate the expression for values of \(t\) close to \(0\):

\begin{datatable}{rcccc}
  \(t\) & \(-0.1\) & \(-0.01\) & 0.01 & 0.1 \\
  \midrule
  \(\frac{\sqrt{t^{2} + 9} - 3}{t^{2}}\) & 0.166620 & 0.166666 & 0.166666 & 0.166620 \\
\end{datatable}
These values are close to \(\tfrac{1}{6} = 0.166666\dots\), so we guess \[ \lim_{t \to 0} \frac{\sqrt{t^{2} + 9} - 3}{t^{2}} = \tfrac{1}{6}. \]
\end{envsolution}
\end{workedexample}
\begin{envcaution}{Caution}{1.5}{}
\begin{raggedpara}
A table of values can only ever \emph{suggest} a limit; it can never confirm one, and it can mislead. Consider \(f(x) = \sin(\pi / x)\) as \(x \to 0\). Sampling \(x = 0.1, 0.01, 0.001, \dots\) gives \(\sin(10\pi), \sin(100\pi), \dots = 0\) every time, suggesting the limit is \(0\). But sampling \(x = 0.3, 0.03, 0.003, \dots\) gives \(\sin\!\left(\tfrac{10\pi}{3}\right), \sin\!\left(\tfrac{100\pi}{3}\right), \dots = -\tfrac{\sqrt{3}}{2}\) every time, suggesting \(-\tfrac{\sqrt{3}}{2}\).
\end{raggedpara}
\begin{raggedpara}
Because \(\pi / x\) grows without bound as \(x \to 0\), the sine runs through more and more complete cycles. In truth, near \(0\) the graph keeps oscillating between \(-1\) and \(1\), no matter how closely we zoom in. The function \(f\) runs through every value in \([-1, 1]\) infinitely often as \(x \to 0\), so the limit does not exist. The definition requires \(f(x)\) to be close to a single \(L\) for \emph{all} \(x\) sufficiently near \(0\), not merely for a convenient list of samples.
\end{raggedpara}
\end{envcaution}
We can now answer the question raised at the start of the section. Suppose a ball thrown upward has height \(s(t) = 30t - 5t^{2}\) metres at time \(t\) seconds. Its average velocity over the interval from \(t\) to \(4\) is

\[ \begin{aligned}
    \frac{s(4) - s(t)}{4 - t}
      &= \frac{40 - (30t - 5t^{2})}{4 - t} \\[2pt]
      &= \frac{5(t - 2)(t - 4)}{4 - t} \\[2pt]
      &= \frac{5(t - 2)(t - 4)}{-(t - 4)} \\[2pt]
      &= -5(t - 2), \qquad t \neq 4.
  \end{aligned} \]

The \emph{instantaneous} velocity at \(t = 4\), written \(v(4)\), is the value this average approaches as the interval shrinks to nothing:

\[ v(4) = \lim_{t \to 4} \frac{s(4) - s(t)}{4 - t} = \lim_{t \to 4} \bigl[-5(t - 2)\bigr] = -10 \text{ m/s}. \]

Notice that the original ratio has the indeterminate form \(\tfrac{0}{0}\) at \(t = 4\) (a form taken up systematically in §1.5): both \(s(4) - s(t)\) and \(4 - t\) tend to \(0\). The cancellation that produced \(-5(t - 2)\) was what made the limit accessible.

The limits above were guessed by looking at function values near the point of interest. In the following sections, we will develop systematic methods for evaluating limits and study more subtle behavior such as one-sided limits and infinite limits.


\sechead{1.4}{One-Sided and Infinite Limits}
Near a point \(a\), a function can behave differently depending on which side \(x\) approaches from, and in some cases the values of the function do not approach any finite number at all. These two observations lead to one-sided limits and infinite limits.

\subsechead{One-sided limits}
Sometimes we only want to consider one side of the approach to \(a\) — for instance because the function is defined by different formulas on the two sides, or because the graph jumps at \(a\).

\begin{envdefinition}{Definition}{1.10}{}
We write \[ \lim_{x \to a^{-}} f(x) = L \] and say that the \emph{left-hand limit} of \(f(x)\) as \(x\) approaches \(a\) is \(L\) if \(f(x)\) approaches \(L\) as \(x\) approaches \(a\) through values less than \(a\). Similarly, \[ \lim_{x \to a^{+}} f(x) = L \] is the \emph{right-hand limit}, describing the approach through values greater than \(a\).
\end{envdefinition}
\begin{figureblock}
  \includegraphics[width=90.49mm]{ch01/figs/one-sided-limits}
\figcaption{Figure 1.14}{A function whose left-hand and right-hand limits at \(a\) are different.}
\end{figureblock}
The two one-sided limits together determine the two-sided limit.

\begin{envproposition}{Proposition}{1.5}{Criterion for a two-sided limit}
\[ \lim_{x \to a} f(x) = L \quad \Longleftrightarrow \quad \lim_{x \to a^{-}} f(x) = L \ \text{ and } \ \lim_{x \to a^{+}} f(x) = L. \]
\end{envproposition}
If the two one-sided limits differ, then the two-sided limit does not exist.

\begin{workedexample}
\begin{envexample}{Example}{1.19}{}
Suppose \(f\) is defined on all of \(\mathbb{R}\).

\begin{enumerate}
  \item If \(\lim_{x \to -2} f(x) = 16\), what is \(\lim_{x \to -2^{-}} f(x)\)?
  \item If instead all we know is \(\lim_{x \to -2^{-}} f(x) = 16\), what can we say about \(\lim_{x \to -2} f(x)\)?
\end{enumerate}
\end{envexample}
\begin{envsolution}{Solution}{}{}
\begin{enumerate}
  \item For the two-sided limit to equal \(16\), \emph{both} one-sided limits must equal \(16\); in particular \(\lim_{x \to -2^{-}} f(x) = 16\).
  \item Not enough information. If the right-hand limit also equals \(16\), then \(\lim_{x \to -2} f(x) = 16\); if it differs or fails to exist, the two-sided limit does not exist.
\end{enumerate}
The criterion transfers information from the two-sided limit to each side, but only \emph{both} sides together determine the two-sided limit.
\end{envsolution}
\end{workedexample}
\begin{workedexample}
\begin{envexample}{Example}{1.20}{}
Let \[ f(x) = \begin{cases} x + 1, & x < 2,\\[4pt] 4 - x, & x \ge 2. \end{cases} \] Evaluate \(\lim_{x \to 2^{-}} f(x)\) and \(\lim_{x \to 2^{+}} f(x)\), and determine whether \(\lim_{x \to 2} f(x)\) exists.
\end{envexample}
\begin{envsolution}{Solution}{}{}
From the left we use the branch \(f(x) = x + 1\), so \(\lim_{x \to 2^{-}} f(x) = 3\).

From the right we use the branch \(f(x) = 4 - x\), so \(\lim_{x \to 2^{+}} f(x) = 2\).

Since the left-hand and right-hand limits differ, \(\lim_{x \to 2} f(x)\) does not exist.

\begin{figureblock}
  \includegraphics[width=95.78mm]{ch01/figs/piecewise-jump}
\figcaption{Figure 1.15}{The jump at \(x = 2\) in Example 1.20.}
\end{figureblock}
\end{envsolution}
\end{workedexample}
\subsechead{Infinite limits}
In some cases, the values of \(f(x)\) do not approach any finite number as \(x \to a\). Instead, they grow arbitrarily large positive or arbitrarily large negative.

\begin{workedexample}
\begin{envexample}{Example}{1.21}{}
Consider \(f(x) = \frac{1}{x^{2}}\). What happens to \(f(x)\) as \(x \to 0\)?
\end{envexample}
\begin{envsolution}{Solution}{}{}
If \(x\) is close to \(0\), then \(x^{2}\) is a very small positive number, so \(\frac{1}{x^{2}}\) is very large. For instance, \[ \frac{1}{(0.1)^{2}} = 100, \qquad \frac{1}{(0.01)^{2}} = 10{,}000. \] The values of \(\frac{1}{x^{2}}\) grow without bound as \(x \to 0\), so no finite \(L\) satisfies the informal definition of the limit; the limit does not exist.
\end{envsolution}
\end{workedexample}
\begin{workedexample}
\begin{envexample}{Example}{1.22}{}
Evaluate \[ \lim_{x \to 0^{-}} \frac{1}{x}, \qquad \lim_{x \to 0^{+}} \frac{1}{x}, \qquad\text{and}\qquad \lim_{x \to 0} \frac{1}{x}, \] and contrast the results with \(\lim_{x \to 0} 1/x^{2}\).
\end{envexample}
\begin{envsolution}{Solution}{}{}
As \(x\) approaches \(0\) from the left, \(1/x\) is negative with larger and larger magnitude, so \[ \lim_{x \to 0^{-}} \frac{1}{x} = -\infty. \] (The limit-value notation \(= -\infty\) is made precise in Definition 1.11 below.) From the right, \(1/x\) is positive and grows without bound, so \[ \lim_{x \to 0^{+}} \frac{1}{x} = \infty. \] Because the two one-sided behaviors disagree, the two-sided limit does not exist. Since the values do not all grow in the \emph{same} signed direction, we do \runin{not} write \(\lim_{x \to 0} \frac{1}{x} = \infty\). Contrast \(1/x^{2}\): both sides grow positively, so there we do write \(\lim_{x \to 0} \frac{1}{x^{2}} = \infty\) (Figure 1.16).

\begin{figureblock}
  \includegraphics[width=53.45mm]{ch01/figs/recip-x-vs-x2-1}\hspace{6mm}\includegraphics[width=53.45mm]{ch01/figs/recip-x-vs-x2-2}
\figcaption{Figure 1.16}{Comparing \(1/x^{2}\) and \(1/x\) near \(0\).}
\end{figureblock}
\end{envsolution}
\end{workedexample}
To capture this behavior, we introduce the notation \[ \lim_{x \to 0} \frac{1}{x^{2}} = \infty. \] The symbol \(\infty\) does not represent a real number; it records that the function’s values grow without bound.

\begin{envdefinition}{Definition}{1.11}{}
Let \(f\) be defined on both sides of \(a\), except possibly at \(a\) itself. We write \[ \lim_{x \to a} f(x) = \infty \] if the values of \(f(x)\) can be made arbitrarily large by taking \(x\) sufficiently close to \(a\), with \(x \neq a\). We call such a limit an \emph{infinite limit}. Similarly, \[ \lim_{x \to a} f(x) = -\infty \] if the values of \(f(x)\) can be made arbitrarily large negative by taking \(x\) sufficiently close to \(a\), with \(x \neq a\).
\end{envdefinition}
\begin{envcaution}{Caution}{1.6}{}
Writing \(\lim_{x \to a} f(x) = \infty\) does \emph{not} mean that the limit exists in the ordinary sense. It is shorthand recording a particular way the limit \emph{fails} to exist as a finite real number: the values of \(f\) grow without bound. The symbol \(\infty\) cannot be added, multiplied, or otherwise manipulated like an ordinary number. Reading “\(= \infty\)” as though it were an ordinary limit value is a common slip.
\end{envcaution}
One-sided versions, \[ \begin{aligned} \lim_{x \to a^{-}} f(x) &= \infty, &\qquad \lim_{x \to a^{+}} f(x) &= \infty, \\ \lim_{x \to a^{-}} f(x) &= -\infty, &\qquad \lim_{x \to a^{+}} f(x) &= -\infty, \end{aligned} \] are defined analogously. The four cases are illustrated in Figure 1.17.

\begin{figureblock}
  \begin{minipage}{\linewidth}\centering
  \includegraphics[width=64.03mm]{ch01/figs/one-sided-infinite-1}\hspace{6mm}\includegraphics[width=64.03mm]{ch01/figs/one-sided-infinite-2}\\[6pt]
  \includegraphics[width=64.03mm]{ch01/figs/one-sided-infinite-3}\hspace{6mm}\includegraphics[width=64.03mm]{ch01/figs/one-sided-infinite-4}
  \figcaption{Figure 1.17}{The four typical one-sided infinite limits.}
  \end{minipage}
\end{figureblock}
\subsechead{Vertical asymptotes}
Infinite limits give a precise way to describe vertical asymptotes — the lines a graph shoots up or down along.

\begin{envdefinition}{Definition}{1.12}{}
The vertical line \(x = a\) is called a \emph{vertical asymptote} of the curve \(y = f(x)\) if at least one of the following holds: \[ \lim_{x \to a^{-}} f(x) = \pm\infty \qquad \text{or} \qquad \lim_{x \to a^{+}} f(x) = \pm\infty. \]
\end{envdefinition}
\begin{envstrategy}{Strategy}{1.3}{Locating vertical asymptotes}
\begin{steps}
  \item Find the candidate points: places where the formula is undefined, the graph breaks, or the expression otherwise looks suspect — though still approachable from within the domain. For a rational function, these are usually the zeros of the denominator.
  \item At each candidate \(x = a\), examine the one-sided limits. When the expression has the form \(\frac{\text{nonzero}}{0}\), determine the sign of the small denominator on each side to decide between \(+\infty\) and \(-\infty\).
  \item If at least one one-sided limit is \(+\infty\) or \(-\infty\), the line \(x = a\) is a vertical asymptote. If instead the problematic factor cancels completely and the resulting one-sided limits are finite, then there is no vertical asymptote there.
\end{steps}
\end{envstrategy}
\begin{workedexample}
\begin{envexample}{Example}{1.23}{}
Does the curve \(y = \frac{2x}{x - 3}\) have a vertical asymptote?
\end{envexample}
\begin{envsolution}{Solution}{}{}
The denominator is zero only at \(x = 3\), so \(x = 3\) is the only candidate. As \(x \to 3^{+}\), the denominator \(x - 3\) is a small positive number while the numerator \(2x\) is close to \(6\), so \[ \lim_{x \to 3^{+}} \frac{2x}{x - 3} = \infty. \] As \(x \to 3^{-}\), the denominator is a small negative number while the numerator is still close to \(6\), so \[ \lim_{x \to 3^{-}} \frac{2x}{x - 3} = -\infty. \] At least one one-sided limit is infinite, so the line \(x = 3\) is a vertical asymptote, as shown in Figure 1.18.

\begin{figureblock}
  \includegraphics[width=95.78mm]{ch01/figs/vertical-asymptote}
\figcaption{Figure 1.18}{The vertical asymptote at \(x = 3\).}
\end{figureblock}
\end{envsolution}
\end{workedexample}
\begin{workedexample}
\begin{envexample}{Example}{1.24}{}
Suppose \(f(x) = \frac{g(x)}{x^{2} - 9}\) for some function \(g\). Must the line \(x = -3\) be a vertical asymptote of \(y = f(x)\)?
\end{envexample}
\begin{envsolution}{Solution}{}{}
No. The denominator does vanish at \(x = -3\), so \(x = -3\) is a \emph{candidate} — but whether it is an asymptote depends on the numerator. If \(g(x) = 1\), then \(f(x) = \frac{1}{x^{2} - 9}\) blows up near \(-3\), and \(x = -3\) is a vertical asymptote. But if \(g(x) = x^{2} - 9\), then \(f(x) = 1\) at every point of its domain, and there is no asymptote at all. A zero denominator marks a place to \emph{investigate} with one-sided limits, not a guaranteed asymptote.
\end{envsolution}
\end{workedexample}
\begin{workedexample}
\begin{envexample}{Example}{1.25}{}
Guess the value of \(\lim_{x \to 0^{+}} \ln x\).
\end{envexample}
\begin{envsolution}{Solution}{}{}
The natural logarithm \(\ln x\) is the logarithm to the base \(e \approx 2.718\); it is studied thoroughly in Chapter 4, and for now the graph in Figure 1.19 supplies everything we need. Since \(\ln x\) is defined only for \(x > 0\), the point \(x = 0\) sits at the boundary of the domain \((0, \infty)\) and can only be approached from within it — that is, from the right. As \(x \to 0^{+}\), the values of \(\ln x\) decrease without bound, so \[ \lim_{x \to 0^{+}} \ln x = -\infty. \] Therefore \(x = 0\) is a vertical asymptote of the graph of \(y = \ln x\).

\begin{figureblock}
  \includegraphics[width=90.49mm]{ch01/figs/ln-asymptote}
\figcaption{Figure 1.19}{The vertical asymptote of \(y = \ln x\).}
\end{figureblock}
\end{envsolution}
\end{workedexample}

\sechead{1.5}{Limit Laws and Computational Techniques}
In the previous sections we used tables and graphs to guess the values of limits. We now record algebraic properties that let us compute limits more systematically. These properties can be justified rigorously once we have the precise definition of a limit (see the next section); for now, they provide a reliable tool-kit.

\subsechead{The limit laws}
Each of these laws says something we already expect. If \(f(x)\) is settling toward \(L\) and \(g(x)\) toward \(M\) as \(x\) nears \(a\), then \(f(x) + g(x)\) ought to settle toward \(L + M\), the product \(f(x)\,g(x)\) toward \(LM\), and so on. In short, the laws assert that combining the functions first and then taking the limit gives the same result as taking the limits first and then combining.

\begin{envtheorem}{Theorem}{1.2}{Limit laws}
Suppose \(c\) is a constant and both \(\lim_{x \to a} f(x)\) and \(\lim_{x \to a} g(x)\) exist. Then \[ \begin{aligned} \lim_{x \to a}\bigl(f(x) + g(x)\bigr) &= \lim_{x \to a} f(x) + \lim_{x \to a} g(x),\\ \lim_{x \to a}\bigl(f(x) - g(x)\bigr) &= \lim_{x \to a} f(x) - \lim_{x \to a} g(x),\\ \lim_{x \to a}\bigl(c \, f(x)\bigr) &= c \lim_{x \to a} f(x),\\ \lim_{x \to a}\bigl(f(x) g(x)\bigr) &= \left(\lim_{x \to a} f(x)\right)\!\left(\lim_{x \to a} g(x)\right). \end{aligned} \] If \(\lim_{x \to a} g(x) \neq 0\), then \[ \lim_{x \to a} \frac{f(x)}{g(x)} = \frac{\lim_{x \to a} f(x)}{\lim_{x \to a} g(x)}. \] Moreover, for every positive integer \(n\), \[ \lim_{x \to a} \bigl[f(x)\bigr]^{n} = \left[\lim_{x \to a} f(x)\right]^{n}. \]
\end{envtheorem}
In words, the limit of a sum / difference / constant multiple / product / integer power is the corresponding combination of the limits, and the limit of a quotient is the quotient of the limits whenever the denominator’s limit is not zero.

\begin{envcaution}{Caution}{1.7}{}
The quotient law applies only when the denominator’s limit is \emph{nonzero}. If \(\lim_{x \to a} g(x) = 0\), the law says nothing, and you may not conclude that the quotient’s limit fails to exist or that it takes any particular value. A resulting \(\tfrac{0}{0}\) is not a verdict about the limit. It is a signal to do more work, usually algebraic simplification, as in the examples below.
\end{envcaution}
The laws apply even when we have no formula for the functions involved, as the next example shows.

\begin{workedexample}
\begin{envexample}{Example}{1.26}{}
Suppose \(\lim_{x \to -1} f(x) = -1\). Evaluate \[ \lim_{x \to -1} \frac{x\,f(x) + 3}{2 f(x) + 1}. \]
\end{envexample}
\begin{envsolution}{Solution}{}{}
We never need a formula for \(f\). By the product law, \(\lim_{x \to -1} x f(x) = (-1)(-1) = 1\); by the sum and constant-multiple laws the numerator tends to \(1 + 3 = 4\) and the denominator tends to \(2(-1) + 1 = -1\). The denominator’s limit is nonzero, so the quotient law applies: \[ \lim_{x \to -1} \frac{x f(x) + 3}{2 f(x) + 1} = \frac{4}{-1} = -4. \]
\end{envsolution}
\end{workedexample}
\subsechead{Direct substitution}
A first consequence of the limit laws is that polynomials and rational functions behave in the easiest way possible.

\begin{envproposition}{Proposition}{1.6}{}
For every positive integer \(n\), \[ \begin{aligned} \lim_{x \to a} x &= a,\\ \lim_{x \to a} x^{n} &= a^{n}. \end{aligned} \] Consequently, if \(p(x) = a_{n} x^{n} + a_{n-1} x^{n - 1} + \dots + a_{0}\) is a polynomial, then \[ \lim_{x \to a} p(x) = p(a). \] If \(r(x)\) is a rational function and \(r(a)\) is defined, then \[ \lim_{x \to a} r(x) = r(a). \]
\end{envproposition}
For polynomials and rational functions where the denominator does not vanish at \(a\), direct substitution is therefore the fastest route to the limit.

\begin{workedexample}
\begin{envexample}{Example}{1.27}{}
Evaluate \[ \lim_{x \to 1} \frac{x^{3} + 2x^{2} - 1}{5 - 3x}. \]
\end{envexample}
\begin{envsolution}{Solution}{}{}
The numerator and denominator are polynomials, so their limits at \(x = 1\) are computed by direct substitution: \[ \begin{aligned} \lim_{x \to 1}\bigl(x^{3} + 2x^{2} - 1\bigr) &= 1 + 2 - 1 = 2,\\ \lim_{x \to 1}(5 - 3x) &= 5 - 3 = 2. \end{aligned} \] The denominator’s limit is nonzero, so the quotient law gives \[ \lim_{x \to 1} \frac{x^{3} + 2x^{2} - 1}{5 - 3x} = \frac{2}{2} = 1. \]
\end{envsolution}
\end{workedexample}
\subsechead{Algebraic simplification}
When direct substitution produces an indeterminate form such as \(\tfrac{0}{0}\), we first simplify the expression and then take the limit.

\begin{workedexample}
\begin{envexample}{Example}{1.28}{}
\begin{enumerate}
  \item Find functions \(f\) and \(g\) with \(\lim_{x \to 3} f(x) = \lim_{x \to 3} g(x) = 0\) and \(\lim_{x \to 3} \frac{f(x)}{g(x)} = 10\).
  \item What are the possible values of \(\lim_{x \to 3} \frac{f(x)}{g(x)}\), given only that both \(f\) and \(g\) tend to \(0\)?
\end{enumerate}
\end{envexample}
\begin{envsolution}{Solution}{}{}
\begin{enumerate}
  \item Take \(f(x) = 10(x - 3)\) and \(g(x) = x - 3\). Both tend to \(0\) at \(3\), yet \(\frac{f(x)}{g(x)} = 10\) wherever it is defined, so the quotient’s limit is \(10\).
  \item Replacing \(10\) by any constant \(c\) gives limit \(c\). The choice \(f(x) = x - 3\), \(g(x) = (x - 3)^{3}\) gives \(\infty\), since the quotient simplifies to \(\frac{1}{(x - 3)^{2}}\); reversing the sign of \(g\) turns the quotient into \(-\frac{1}{(x - 3)^{2}}\), which gives \(-\infty\). Finally, \(f(x) = x - 3\), \(g(x) = |x - 3|\) makes the limit fail to exist: this quotient is \(1\) for \(x > 3\) but \(-1\) for \(x < 3\), so the two one-sided limits disagree.
\end{enumerate}
Therefore, when substitution gives \(\tfrac{0}{0}\), the limit may be \emph{any} number, may be infinite, or may fail to exist. That is exactly why the form is called indeterminate, and why we must simplify before judging.
\end{envsolution}
\end{workedexample}
\begin{workedexample}
\begin{envexample}{Example}{1.29}{}
Evaluate \[ \lim_{x \to 1} \frac{x - 1}{x^{2} - 1}. \]
\end{envexample}
\begin{envsolution}{Solution}{}{}
Direct substitution gives \(\tfrac{0}{0}\), so we simplify: \[ \begin{aligned} \frac{x - 1}{x^{2} - 1} &= \frac{x - 1}{(x - 1)(x + 1)}\\ &= \frac{1}{x + 1}, \qquad x \neq 1. \end{aligned} \] Therefore \[ \lim_{x \to 1} \frac{x - 1}{x^{2} - 1} = \lim_{x \to 1} \frac{1}{x + 1} = \tfrac{1}{2}. \]

\begin{figureblock}
  \includegraphics[width=85.19mm]{ch01/figs/removable-hole}
\figcaption{Figure 1.20}{The removable hole at \(x = 1\); its height is \(\tfrac{1}{2}\).}
\end{figureblock}
\end{envsolution}
\end{workedexample}
\begin{workedexample}
\begin{envexample}{Example}{1.30}{}
Evaluate \[ \lim_{t \to 0} \frac{\sqrt{t^{2} + 9} - 3}{t^{2}}. \]
\end{envexample}
\begin{envsolution}{Solution}{}{}
Direct substitution gives \(\tfrac{0}{0}\), so we rationalize the numerator: \[ \begin{aligned} \frac{\sqrt{t^{2} + 9} - 3}{t^{2}} \cdot \frac{\sqrt{t^{2} + 9} + 3}{\sqrt{t^{2} + 9} + 3} &= \frac{t^{2} + 9 - 9}{t^{2}\bigl(\sqrt{t^{2} + 9} + 3\bigr)}\\ &= \frac{1}{\sqrt{t^{2} + 9} + 3}, \qquad t \neq 0. \end{aligned} \] Therefore \[ \lim_{t \to 0} \frac{\sqrt{t^{2} + 9} - 3}{t^{2}} = \lim_{t \to 0} \frac{1}{\sqrt{t^{2} + 9} + 3} = \tfrac{1}{6}. \]
\end{envsolution}
\end{workedexample}
\begin{workedexample}
\begin{envexample}{Example}{1.31}{}
Evaluate \[ \lim_{t \to \frac{1}{2}} \frac{\frac{1}{3t^{2}} + \frac{1}{t^{2} - 1}}{2t - 1}. \]
\end{envexample}
\begin{envsolution}{Solution}{}{}
Direct substitution gives \(\tfrac{0}{0}\). Combining the two fractions in the numerator over the common denominator \(3t^{2}(t^{2} - 1)\), \[ \begin{aligned} \frac{1}{3t^{2}} + \frac{1}{t^{2} - 1} &= \frac{(t^{2} - 1) + 3t^{2}}{3t^{2}(t^{2} - 1)} = \frac{4t^{2} - 1}{3t^{2}(t^{2} - 1)} \\ &= \frac{(2t - 1)(2t + 1)}{3t^{2}(t^{2} - 1)}. \end{aligned} \] Dividing by \(2t - 1\) cancels the troublesome factor, so \[ \begin{aligned} \lim_{t \to \frac{1}{2}} \frac{\frac{1}{3t^{2}} + \frac{1}{t^{2} - 1}}{2t - 1} &= \lim_{t \to \frac{1}{2}} \frac{2t + 1}{3t^{2}(t^{2} - 1)} \\ &= \frac{2}{\tfrac{3}{4}\bigl(-\tfrac{3}{4}\bigr)} = -\frac{32}{9}. \end{aligned} \]
\end{envsolution}
\end{workedexample}
\subsechead{Limits of piecewise-defined functions}
When a function is defined by different formulas on either side of \(a\), the one-sided limits can disagree, so we analyse each branch separately.

\begin{workedexample}
\begin{envexample}{Example}{1.32}{}
Evaluate \(\lim_{x \to 0} |x|\).
\end{envexample}
\begin{envsolution}{Solution}{}{}
Using the piecewise description \[ |x| = \begin{cases} x, & x \ge 0,\\[4pt] -x, & x < 0, \end{cases} \] we compute \[ \begin{aligned} \lim_{x \to 0^{+}} |x| &= \lim_{x \to 0^{+}} x = 0, \\ \lim_{x \to 0^{-}} |x| &= \lim_{x \to 0^{-}} (-x) = 0. \end{aligned} \] Both one-sided limits equal \(0\), so \(\lim_{x \to 0} |x| = 0\).
\end{envsolution}
\end{workedexample}
\begin{workedexample}
\begin{envexample}{Example}{1.33}{}
The greatest integer function \([x]\), often also written \(\lfloor x\rfloor\), is the largest integer less than or equal to \(x\), so that \([4] = 4\), \([4.1] = 4\), and \([4.8] = 4\). Determine whether \(\lim_{x \to 3} [x]\) exists.
\end{envexample}
\begin{envsolution}{Solution}{}{}
As \(x \to 3^{+}\), the values of \(x\) lie just above \(3\), so \([x] = 3\) and \(\lim_{x \to 3^{+}} [x] = 3\). As \(x \to 3^{-}\), the values of \(x\) lie just below \(3\), so \([x] = 2\) and \(\lim_{x \to 3^{-}} [x] = 2\). The two one-sided limits differ, so \(\lim_{x \to 3} [x]\) does not exist.

\begin{figureblock}
  \includegraphics[width=93.13mm]{ch01/figs/floor-function}
\figcaption{Figure 1.21}{The greatest integer function \([x]\) has a jump at every integer.}
\end{figureblock}
\end{envsolution}
\end{workedexample}
\subsechead{The squeeze theorem}
The limit laws require that both constituent limits exist, but sometimes one of the factors has no limit at all. A useful technique is to trap the expression between two simpler functions whose limits agree.

\begin{envtheorem}{Theorem}{1.3}{Squeeze theorem}
Suppose that \[ g(x) \le f(x) \le h(x) \] for all \(x\) near \(a\), except possibly at \(x = a\). If \[ \lim_{x \to a} g(x) = L, \qquad \lim_{x \to a} h(x) = L, \] then \(\lim_{x \to a} f(x) = L\).
\end{envtheorem}
\begin{figureblock}
  \includegraphics[width=111.65mm]{ch01/figs/squeeze-theorem}
\figcaption{Figure 1.22}{The squeeze theorem: \(f\) is trapped between \(g\) and \(h\), which share the same limit \(L\).}
\end{figureblock}
\begin{workedexample}
\begin{envexample}{Example}{1.34}{}
Evaluate \[ \lim_{x \to 0} x^{2} \sin\!\left(\frac{1}{x}\right). \]
\end{envexample}
\begin{envsolution}{Solution}{}{}
We cannot apply the product law because \(\lim_{x \to 0} \sin(1/x)\) does not exist. Instead, we use the bound \[ -1 \le \sin\!\left(\frac{1}{x}\right) \le 1 \] valid for all \(x \neq 0\). Multiplying by \(x^{2} \ge 0\), \[ -x^{2} \le x^{2} \sin\!\left(\frac{1}{x}\right) \le x^{2}. \] Since \(\lim_{x \to 0}(-x^{2}) = 0 = \lim_{x \to 0} x^{2}\), the squeeze theorem yields \[ \lim_{x \to 0} x^{2} \sin\!\left(\frac{1}{x}\right) = 0. \]

\begin{figureblock}
  \includegraphics[width=106.36mm]{ch01/figs/squeeze-x2sin}
\figcaption{Figure 1.23}{The graph of \(y = x^{2}\sin(1/x)\) squeezed between \(y = -x^{2}\) and \(y = x^{2}\).}
\end{figureblock}
\end{envsolution}
\end{workedexample}
\begin{workedexample}
\begin{envexample}{Example}{1.35}{}
Evaluate \[ \lim_{x \to 0} x \cos\!\left(\frac{1}{x}\right). \]
\end{envexample}
\begin{envsolution}{Solution}{}{}
Since \(\left|\cos(1/x)\right| \le 1\) for all \(x \neq 0\), \[ \left|x \cos\!\left(\frac{1}{x}\right)\right| \le |x|, \] equivalently \[ -|x| \le x \cos\!\left(\frac{1}{x}\right) \le |x|. \] Since \(\lim_{x \to 0}(-|x|) = 0 = \lim_{x \to 0} |x|\), the squeeze theorem yields \[ \lim_{x \to 0} x \cos\!\left(\frac{1}{x}\right) = 0. \]
\end{envsolution}
\end{workedexample}
\begin{envremark}{Remark}{1.5}{}
The squeeze theorem also plays a central role beyond routine exercises. It is the standard tool for proving the fundamental trigonometric limit \(\lim_{x \to 0} \frac{\sin x}{x} = 1\), where the ratio is trapped between geometric area bounds. That limit, in turn, is what makes the derivative formulas for \(\sin x\) and \(\cos x\) work in Chapter 3.
\end{envremark}
\begin{envstrategy}{Strategy}{1.4}{Computing a limit \(\lim_{x \to a} f(x)\)}
\begin{steps}
  \item \emph{Try direct substitution first.} If \(f(a)\) is defined and the expression is a polynomial or rational function with a nonzero denominator at \(a\), then the limit equals \(f(a)\).
  \item \emph{If substitution gives an indeterminate form} such as \(\tfrac{0}{0}\), simplify algebraically: factor and cancel, rationalize a numerator or denominator, or combine fractions.
  \item \emph{For a piecewise function}, evaluate the two one-sided limits on the appropriate branches and check whether they agree.
  \item \emph{If the expression is trapped between two simpler ones} whose limits at \(a\) are known and equal, apply the squeeze theorem.
  \item \emph{If the values grow without bound}, the limit is \(\pm\infty\) and the line \(x = a\) is a vertical asymptote.
\end{steps}
\end{envstrategy}
\begin{workedexample}
\begin{envexample}{Example}{1.36}{}
Find the value of the constant \(a\) for which \[ \lim_{x \to -2} \frac{x^{2} + a x + 3}{x^{2} + x - 2} \] exists, and evaluate the limit for that value of \(a\).
\end{envexample}
\begin{envsolution}{Solution}{}{}
The denominator factors as \((x + 2)(x - 1)\) and tends to \(0\) as \(x \to -2\). If the numerator tended to a nonzero number, the quotient would blow up and the limit could not exist; so the numerator must also tend to \(0\): \(4 - 2a + 3 = 0\), giving \(a = \tfrac{7}{2}\). With this \(a\), \[ \begin{aligned} \lim_{x \to -2} \frac{x^{2} + \tfrac{7}{2}x + 3}{x^{2} + x - 2} &= \lim_{x \to -2} \frac{(x + 2)\bigl(x + \tfrac{3}{2}\bigr)}{(x + 2)(x - 1)} \\ &= \lim_{x \to -2} \frac{x + \tfrac{3}{2}}{x - 1} = \frac{-\tfrac{1}{2}}{-3} = \frac{1}{6}. \end{aligned} \] Notice the flow of reasoning: the \emph{existence} of the limit forced the value of \(a\), and then the techniques of this section finished the computation.
\end{envsolution}
\end{workedexample}

\sechead{1.6}{The Precise Definition of a Limit}
\emph{First reading: this section turns the informal “close to” of §1.3–§1.5 into a precise definition. Chapter 2 uses the limit laws themselves, not the \(\varepsilon\)-\(\delta\) arguments behind them, so you may go on to derivatives after §1.5 and return here when a proof calls for the precise definition. The chain-rule proof of §3.2 is the first that does.}

From the time of Newton and Leibniz, calculus was remarkably successful at solving problems in geometry, mechanics, and astronomy. For a long time, however, there was also controversy about its logical foundations. Early arguments often relied on infinitesimals — quantities thought of as being extremely small but not described in a fully precise way — and mathematicians continued to ask what it really means to say that \(x\) approaches \(a\) or that \(f(x)\) approaches \(L\). The modern resolution, developed in the nineteenth century, is the \(\varepsilon\)-\(\delta\) definition of a limit (stated below in terms of two positive tolerances, \(\varepsilon\) and \(\delta\)), which gave calculus the level of rigor used in modern mathematics.

\begin{envremark}{Remark}{1.6}{Historical note}
The drive to replace intuitive arguments with arithmetic precision is known as the \emph{arithmetization of analysis}. Bolzano anticipated a more rigorous, inequality-based treatment of continuity and convergence, while Cauchy gave limits and continuity a more systematic analytic form. Weierstrass, lecturing at Berlin from the 1860s onward, refined these ideas into the \(\varepsilon\text{-}\delta\) formulation that is essentially the one presented here.
\end{envremark}
In the previous sections we used graphs, tables, and algebraic techniques to study limits, and we stated the limit laws without proof. Those arguments rested on informal phrases such as “\(f(x)\) is close to \(L\)” and “\(x\) is close to \(a\).” To make the arguments rigorous — and, eventually, to prove the limit laws — we need a definition that replaces “close” with numerical precision.

It helps to read \(\varepsilon\) and \(\delta\) as \emph{tolerances}. In engineering and manufacturing one specifies how far a finished quantity may stray from its target: an output tolerance \(\varepsilon\) around the desired value \(L\). One then asks how tightly the input must be controlled near \(a\) to keep the result within that tolerance. The answer is an input tolerance \(\delta\). The definition below makes this everyday promise precise: control the input closely enough, and the output stays within specification.

To motivate the definition, consider the function \[ f(x) = \begin{cases} 2x - 1, & x \neq 3,\\[4pt] 6, & x = 3. \end{cases} \] When \(x\) is close to \(3\) but \(x \neq 3\), the value \(f(x) = 2x - 1\) is close to \(5\). Intuitively, we expect \(\lim_{x \to 3} f(x) = 5\). But how close, exactly? We begin by making the informal phrase “close” numerically precise.

\begin{workedexample}
\begin{envexample}{Example}{1.37}{}
For the function above, how close must \(x\) be to \(3\) so that \(f(x)\) differs from \(5\) by less than \(0.1\)?
\end{envexample}
\begin{envsolution}{Solution}{}{}
We want \(|f(x) - 5| < 0.1\). For \(x \neq 3\), \[ \begin{aligned} |f(x) - 5| &= |(2x - 1) - 5|\\ &= |2x - 6|\\ &= 2 |x - 3|. \end{aligned} \] Requiring \(2 |x - 3| < 0.1\) is equivalent to \(|x - 3| < 0.05\). Therefore, if \(0 < |x - 3| < 0.05\), then \(|f(x) - 5| < 0.1\).
\end{envsolution}
\end{workedexample}
Tightening the tolerance from \(0.1\) to \(0.01\), the same reasoning gives \(|f(x) - 5| < 0.01\) whenever \(0 < |x - 3| < 0.005\). As the tolerance shrinks, so must the neighborhood of \(3\). More generally, for any \(\varepsilon > 0\), the same manipulation gives \(|f(x) - 5| < \varepsilon\) whenever \(0 < |x - 3| < \tfrac{\varepsilon}{2}\); we may therefore take \(\delta = \tfrac{\varepsilon}{2}\). The word \emph{arbitrarily} in the informal definition means precisely: we can answer the question for \emph{every} \(\varepsilon > 0\), not just for a few specific tolerances.

\begin{envdefinition}{Definition}{1.13}{}
Let \(f\) be a function defined on an open interval containing \(a\), except possibly at \(a\) itself. We say that \[ \lim_{x \to a} f(x) = L \] if for every \(\varepsilon > 0\) there exists a \(\delta > 0\) such that \[ |f(x) - L| < \varepsilon \qquad \text{whenever} \qquad 0 < |x - a| < \delta. \]

\begin{informal}
Informally, we can force \(f(x)\) as close to \(L\) as we like by taking \(x\) close enough to \(a\) (but not equal to \(a\)).
\end{informal}
\end{envdefinition}
Figure 1.24 gives a geometric reading of the definition. If \(x\) is restricted to the punctured interval \((a - \delta, a + \delta)\) with \(x \neq a\), then the corresponding values of \(f(x)\) lie in the target strip \((L - \varepsilon, L + \varepsilon)\). The value \(f(a)\) itself may be different from \(L\): the limit depends only on the behavior of \(f(x)\) when \(x\) is close to \(a\), not on what happens at \(x = a\).

\begin{figureblock}
  \includegraphics[width=95.78mm]{ch01/figs/precise-limit}
\figcaption{Figure 1.24}{A punctured \(\delta\)-neighborhood maps into an \(\varepsilon\)-strip around \(L\).}
\end{figureblock}
The condition \(0 < |x - a| < \delta\) says that \(x\) lies within distance \(\delta\) of \(a\), but \(x \neq a\). The inequality \(|f(x) - L| < \varepsilon\) says that \(f(x)\) lies within distance \(\varepsilon\) of \(L\). The key phrase is “for every \(\varepsilon > 0\)”: no matter how small a tolerance \(\varepsilon\) we demand, we must be able to find a corresponding \(\delta\) (Figure 1.25).

\begin{figureblock}
  \includegraphics[width=64.03mm]{ch01/figs/epsilon-delta-dynamic-1}\hspace{6mm}\includegraphics[width=64.03mm]{ch01/figs/epsilon-delta-dynamic-2}
\figcaption{Figure 1.25}{For \(f(x) = 2x - 1\) at \(a = 3\), tightening \(\varepsilon\) forces a narrower \(\delta\)-band.}
\end{figureblock}
It is useful to picture the definition as a challenge-and-response game. A skeptic names a tolerance \(\varepsilon > 0\), stating how close to \(L\) they insist \(f(x)\) must be. Our job is to respond with a \(\delta > 0\) that meets the demand: every \(x\) within \(\delta\) of \(a\) (but not equal to \(a\)) is sent to within \(\varepsilon\) of \(L\). The limit equals \(L\) precisely when we can answer \emph{every} challenge, no matter how small the \(\varepsilon\). This is why \(\varepsilon\) comes first and \(\delta\) second: \(\delta\) is chosen in reply to \(\varepsilon\).

\begin{workedexample}
\begin{envexample}{Example}{1.38}{}
What is wrong with the following attempted definition of \(\lim_{x \to a} f(x) = L\)?

\emph{“Given any \(\delta > 0\), there exists \(\varepsilon > 0\) such that whenever \(|f(x) - L| < \varepsilon\), we have \(0 < |x - a| < \delta\).”}
\end{envexample}
\begin{envsolution}{Solution}{}{}
Two things are reversed. First, the \emph{order of the tolerances}: the challenge must come on the \(y\)-side, “given any \(\varepsilon\),” and only then do we produce a \(\delta\) in response. The attempted version instead lets \(\varepsilon\) be chosen after \(\delta\) is known. Second, the \emph{direction of the implication}: the definition must say that closeness of \(x\) to \(a\) forces closeness of \(f(x)\) to \(L\), that is, \(0 < |x - a| < \delta\) implies \(|f(x) - L| < \varepsilon\), not the other way around. Correcting both gives exactly Definition 1.13.
\end{envsolution}
\end{workedexample}
A first payoff of the \(\varepsilon\)-\(\delta\) definition is that it settles a question the informal picture left open: when a limit exists, its value is unambiguous.

\begin{envproposition}{Proposition}{1.7}{Uniqueness of limits}
If \(\lim_{x \to a} f(x)\) exists, then it is unique.
\end{envproposition}
\begin{envproof}{Proof}{}{}
Suppose that \[ \begin{aligned} \lim_{x \to a} f(x) &= L,\\ \lim_{x \to a} f(x) &= M, \end{aligned} \] and that \(L \neq M\). Let \(\varepsilon = \tfrac{|L - M|}{2}\), half the gap between the two supposed limits. This is the choice that will make the contradiction appear at the end. By the definition of the limit, there exist \(\delta_{1}, \delta_{2} > 0\) such that \[ |f(x) - L| < \varepsilon \quad \text{whenever} \quad 0 < |x - a| < \delta_{1}, \] and \[ |f(x) - M| < \varepsilon \quad \text{whenever} \quad 0 < |x - a| < \delta_{2}. \] If \(0 < |x - a| < \min\{\delta_{1}, \delta_{2}\}\), both inequalities hold. By the triangle inequality \(|p + q| \le |p| + |q|\) (applied with \(p = L - f(x)\) and \(q = f(x) - M\)), \[ \begin{aligned} |L - M| &= |L - f(x) + f(x) - M|\\ &\le |L - f(x)| + |f(x) - M|\\ &< 2\varepsilon\\ &= |L - M|, \end{aligned} \] which is impossible. Therefore \(L = M\).
\end{envproof}
We now verify several limits directly from the definition, building from the simplest case to one that requires an extra step.

\begin{workedexample}
\begin{envexample}{Example}{1.39}{}
Show from the definition that \(\lim_{x \to 2} 5 = 5\).
\end{envexample}
\begin{envsolution}{Solution}{}{}
Given \(\varepsilon > 0\), we need \(|5 - 5| < \varepsilon\) whenever \(0 < |x - 2| < \delta\). But \(|5 - 5| = 0 < \varepsilon\) is \emph{always} true, no matter what \(x\) does. So any choice of \(\delta\) works; we may simply take \(\delta = \varepsilon\) (or \(\delta = 1\), or anything else). This is the one situation where \(\delta\) need not depend on \(\varepsilon\) at all.
\end{envsolution}
\end{workedexample}
\begin{workedexample}
\begin{envexample}{Example}{1.40}{}
Show from the definition that \(\lim_{x \to 3} (4x - 5) = 7\).
\end{envexample}
\begin{envsolution}{Solution}{}{}
Given \(\varepsilon > 0\), we seek \(\delta > 0\) with \[ |(4x - 5) - 7| < \varepsilon \quad \text{whenever} \quad 0 < |x - 3| < \delta. \] Since \[ \begin{aligned} |(4x - 5) - 7| &= |4x - 12|\\ &= 4 |x - 3|, \end{aligned} \] it is enough to require \(4 |x - 3| < \varepsilon\), i.e. \(|x - 3| < \tfrac{\varepsilon}{4}\). We choose \(\delta = \tfrac{\varepsilon}{4}\); then whenever \(0 < |x - 3| < \delta\), \[ \begin{aligned} |(4x - 5) - 7| &= 4 |x - 3|\\ &< 4 \cdot \tfrac{\varepsilon}{4}\\ &= \varepsilon. \end{aligned} \] Therefore \(\lim_{x \to 3} (4x - 5) = 7\).
\end{envsolution}
\end{workedexample}
\begin{workedexample}
\begin{envexample}{Example}{1.41}{}
Show from the definition that \(\lim_{x \to 1} (x^{2} - 5x + 6) = 2\).
\end{envexample}
\begin{envsolution}{Solution}{}{}
Given \(\varepsilon > 0\), we seek \(\delta > 0\) with \[ |(x^{2} - 5x + 6) - 2| < \varepsilon \quad \text{whenever} \quad 0 < |x - 1| < \delta. \] We first simplify: \[ \begin{aligned} |(x^{2} - 5x + 6) - 2| &= |x^{2} - 5x + 4|\\ &= |(x - 1)(x - 4)|. \end{aligned} \] The factor \(|x - 4|\) depends on \(x\). To control it, first assume \(|x - 1| < 1\) (legitimate because we will in the end choose \(\delta \le 1\)). Then \[ \begin{aligned} |x - 1| < 1 &\implies -1 < x - 1 < 1\\ &\implies 0 < x < 2\\ &\implies -4 < x - 4 < -2\\ &\implies |x - 4| < 4. \end{aligned} \] So whenever \(|x - 1| < 1\), we have \(|(x - 1)(x - 4)| < 4 |x - 1|\). To force this below \(\varepsilon\), we also require \(4 |x - 1| < \varepsilon\), i.e. \(|x - 1| < \tfrac{\varepsilon}{4}\). Choosing \[ \delta = \min\!\left\{1, \tfrac{\varepsilon}{4}\right\}, \] we get both inequalities at once, and so \[ \begin{aligned} |(x^{2} - 5x + 6) - 2| &= |(x - 1)(x - 4)|\\ &< 4 |x - 1|\\ &< 4 \cdot \tfrac{\varepsilon}{4}\\ &= \varepsilon. \end{aligned} \] Therefore \(\lim_{x \to 1} (x^{2} - 5x + 6) = 2\).
\end{envsolution}
\end{workedexample}
\begin{envcaution}{Caution}{1.8}{}
The value of \(\delta\) is never unique. Once a valid \(\delta\) has been found, any smaller positive value also works, because shrinking the input neighborhood can only tighten the conclusion. The goal of an \(\varepsilon\text{-}\delta\) proof is to exhibit \emph{one} working \(\delta\), not “the” \(\delta\). In Example 1.40, for instance, \(\delta = \varepsilon/5\) works just as well as \(\delta = \varepsilon/4\).
\end{envcaution}
\begin{workedexample}
\begin{envexample}{Example}{1.42}{}
Using the inequality \(|\sin x| \le |x|\) for all \(x\), show from the definition that \(\lim_{x \to 0} \sin x = 0\).
\end{envexample}
\begin{envsolution}{Solution}{}{}
Given \(\varepsilon > 0\), choose \(\delta = \varepsilon\). If \(0 < |x - 0| < \delta\), then \[ |\sin x - 0| = |\sin x| \le |x| < \delta = \varepsilon. \] Therefore \(\lim_{x \to 0} \sin x = 0\). The given inequality already provides the bound we need, so no factoring and no auxiliary bound are required.
\end{envsolution}
\end{workedexample}
The four examples above share a common pattern, which we record as a general strategy.

\begin{envstrategy}{Strategy}{1.5}{Verifying a limit from the \(\varepsilon\)-\(\delta\) definition}
To show that \(\lim_{x \to a} f(x) = L\):

\begin{steps}
  \item Start with a generic \(\varepsilon > 0\). The task is to find a \(\delta > 0\) (which may depend on \(\varepsilon\)) that makes the definition work.
  \item Work with the quantity \(|f(x) - L|\) and bound it above by an expression involving \(|x - a|\).
  \item Choose \(\delta\) so that the bound above is less than \(\varepsilon\) whenever \(|x - a| < \delta\). If the bound contains an auxiliary factor depending on \(x\), first restrict \(x\) to a small neighborhood of \(a\) so that this factor is bounded. Then take \(\delta\) to be the smaller of the neighborhood radius and the \(\varepsilon\)-based bound.
  \item Conclude that \(0 < |x - a| < \delta\) implies \(|f(x) - L| < \varepsilon\).
\end{steps}
\end{envstrategy}
The \(\varepsilon\)-\(\delta\) machinery also gives precise meaning to the informal definitions of infinite limits from Definition 1.11.

\begin{envdefinition}{Definition}{1.14}{}
We write \[ \lim_{x \to a} f(x) = \infty \] if for every positive number \(M\) there exists \(\delta > 0\) such that \[ f(x) > M \qquad \text{whenever} \qquad 0 < |x - a| < \delta. \] Similarly, \[ \lim_{x \to a} f(x) = -\infty \] if for every negative number \(N\) there exists \(\delta > 0\) such that \[ f(x) < N \qquad \text{whenever} \qquad 0 < |x - a| < \delta. \]
\end{envdefinition}
Intuitively, \(M\) may be any large positive threshold (such as \(100\), \(100^{100}\), or \(1000^{1000}\)): no matter how large the threshold, \(f(x)\) eventually exceeds it once \(x\) is close enough to \(a\). Likewise, \(N\) may be any negative threshold (such as \(-100\), \(-100^{100}\), or \(-1000^{1000}\)), below which \(f(x)\) is eventually forced.

\begin{workedexample}
\begin{envexample}{Example}{1.43}{}
Show from the precise definition that \[ \lim_{x \to 0} \frac{1}{x^{2}} = \infty. \]
\end{envexample}
\begin{envsolution}{Solution}{}{}
Given \(M > 0\), choose \[ \delta = \frac{1}{\sqrt{M}}. \] If \(0 < |x| < \delta\), then \(x^{2} < \delta^{2} = 1/M\), so \[ \frac{1}{x^{2}} > M. \] Since every threshold \(M\) is eventually exceeded once \(x\) is close enough to \(0\), we have \(\lim_{x \to 0} \frac{1}{x^{2}} = \infty\). This is the rigorous counterpart of the informal argument in Example 1.21.
\end{envsolution}
\end{workedexample}
Finally, the \(\varepsilon\)-\(\delta\) definition gives a precise meaning to a property of functions that recurs throughout later chapters. If \(f\) is defined at every point near \(a\), including at \(a\), and if \[ \lim_{x \to a} f(x) = f(a), \] then \(f\) is said to be \emph{continuous at \(a\)}. A systematic study of continuity belongs to a later chapter; we mention the term here because it is expressed directly in terms of limits.

\subsechead{Chapter summary}
Chapter 1 laid down two foundations that look very different from each other. The first was the inverse of a function (§1.1–§1.2). A function can be run backwards exactly when it is \emph{one-to-one}: distinct inputs give distinct outputs. The horizontal line test reads this condition straight off a graph. For such a function the inverse \(f^{-1}\) exists (Theorem 1.1), described exactly by the identities \(f^{-1}(f(x)) = x\) and \(f(f^{-1}(y)) = y\) (Proposition 1.1). Restricting each trigonometric function to a principal interval on which it is one-to-one produced the inverse trigonometric functions \(\arcsin\), \(\arccos\), \(\arctan\) and their three companions, with fixed principal ranges and the reference-triangle method (Strategy 1.2) turning expressions such as \(\cos(\arcsin x)\) into algebra.

The second, and longer, thread was the \emph{limit} (§1.3–§1.6). We began with the informal idea that \(\lim_{x \to a} f(x) = L\) means that \(f(x)\) is as close to \(L\) as we please once \(x\) is close enough to, but not equal to, \(a\). We estimated such limits from tables and graphs (§1.3). One-sided limits refined the picture, with a two-sided limit existing exactly when the two one-sided limits agree (Proposition 1.5); allowing the values to grow without bound gave infinite limits and the vertical asymptotes that mark them (§1.4). The limit laws (Theorem 1.2) then made limits computable: for functions built from polynomials and rational functions, direct substitution works wherever the function is defined (Proposition 1.6), and the indeterminate form \(0/0\) is handled by factoring, rationalization, or the squeeze theorem (Theorem 1.3) (§1.5).

The precise \(\varepsilon\)-\(\delta\) definition (Definition 1.13) finally placed every preceding limit argument on a rigorous footing (§1.6): \(\lim_{x \to a} f(x) = L\) means that for every \(\varepsilon > 0\) there is a \(\delta > 0\) with \(\lvert f(x) - L \rvert < \varepsilon\) whenever \(0 < \lvert x - a \rvert < \delta\), and from it we proved that a limit, when it exists, is unique (Proposition 1.7). The chapter closed by naming continuity, the property that limits make precise. A function \(f\) is continuous at \(a\) when \(\lim_{x \to a} f(x) = f(a)\). Chapter 2 is built on this machinery: the derivative is itself a limit, the limit of secant slopes.
\end{document}
